\section{Eucharistic Hunger, Suffering Within Grace, and the Test of Neighbor-Love}

Friendship with God is never a disembodied affair conducted solely inside private thoughts. Because the Word became flesh and gave himself bodily for the life of the world, charity draws the creature into concrete, visible arenas where love is nourished, tested, and verified: the Eucharistic altar, the crucible of human suffering, and the face of the neighbor.

\subsection*{From Sunday Duty to Eucharistic Hunger}

Justice already establishes why a creature owes God public, communal worship on the Lord's Day. Charity does not abolish that obligation; it transforms the posture with which a believer crosses the church threshold. What begins as faithful obedience to a precept ripens into hunger for the Bread of Life. Christ himself spoke of that meal in the language of fierce longing: ``With desire I have desired to eat this pasch with you, before I suffer'' (Luke 22:15), and ``My flesh is meat indeed: and my blood is drink indeed. He that eateth my flesh and drinketh my blood abideth in me: and I in him'' (John 6:56--57).

In the Tertia Pars of the \emph{Summa Theologiae} (qq.\,73, 79--80), St.\ Thomas Aquinas sets the Eucharist at the summit of the whole sacramental economy. Whereas Baptism is the door of the spiritual life---incorporating a person into Christ through faith---``the Eucharist is, as it were, the consummation of the spiritual life, and the end of all the sacraments,'' rightly named ``the sacrament of Charity, which is `the bond of perfection' (Col.\ 3:14)'' (\emph{ST} III, q.\,73, a.\,3, corpus and ad 3).\endnote{Aquinas, \emph{ST} III, q.\,73, a.\,3, corpus and ad 3. The \emph{Catechism of the Catholic Church} (no.\ 1324) and St.\ John Paul II's encyclical \emph{Ecclesia de Eucharistia} (2003, nos.\ 1, 11) teach this same synthesis, naming the Eucharist the source and summit of the Christian life, uniting the Church to Christ's Paschal Sacrifice.} Every other sacrament prepares for the Eucharist or flows from it, because its ultimate spiritual reality (\emph{res sacramenti}) is the unity of Christ's Mystical Body in charity.

Aquinas draws a sharp contrast between bodily food and this sacramental food (\emph{ST} III, q.\,73, a.\,3, ad 2; q.\,79, a.\,1). When we eat ordinary bread, our body assimilates the bread, breaking it down into our own flesh and blood. In the Eucharist, the direction of assimilation reverses:

\begin{quote}
The difference between corporeal and spiritual food lies in this, that the former is changed into the substance of the person nourished, and consequently it cannot avail for supporting life except it be partaken of; but spiritual food changes man into itself, according to that saying of Augustine (\emph{Confess.}\ vii), that he heard the voice of Christ as it were saying to him: ``Nor shalt thou change Me into thyself, as food of thy flesh, but thou shalt be changed into Me.''\endnote{Aquinas, \emph{ST} III, q.\,73, a.\,3, ad 2. In \emph{ST} III, q.\,79, a.\,1, corpus and ad 1--2, Aquinas shows that the Eucharist sustains, increases, restores, and delights the spiritual life---not only increasing habitual grace and charity, but arousing charity into living act.}
\end{quote}

Yet precisely because the Eucharist signifies and effects real union with Christ and his Church, Catholic doctrine preserves a clear boundary between \emph{subjective desire} and \emph{worthy sacramental reception}. Longing to approach the altar does not by itself erase grave sin. St.\ Paul's warning stands in full force: ``But let a man prove himself: and so let him eat of that bread and drink of the chalice. For he that eateth and drinketh unworthily eateth and drinketh judgment to himself, not discerning the body of the Lord'' (1 Cor 11:28--29).

Aquinas explains why in \emph{ST} III, q.\,80, a.\,4: every sacrament is a sign of a spiritual reality. To receive the Eucharist is to enact a public, bodily sign declaring that one is united with Christ and his Mystical Body through living faith and charity. Therefore, ``whoever receives this sacrament while in mortal sin, is guilty of lying to this sacrament, and consequently of sacrilege, because he profanes the sacrament: and therefore he sins mortally.''\endnote{Aquinas, \emph{ST} III, q.\,80, a.\,4, corpus. See also the Council of Trent, Session XIII, Decree on the Most Holy Eucharist, ch.\ 7; \emph{CCC} 1385 and 1415; and St.\ John Paul II, \emph{Ecclesia de Eucharistia}, nos.\ 36--37: anyone conscious of grave sin must receive sacramental absolution in the Sacrament of Penance before approaching Holy Communion.} To the objection that the Eucharist is spiritual medicine and therefore belongs to the spiritually sick, Aquinas answers with a medical distinction (\emph{ST} III, q.\,80, a.\,4, ad 2): not every medicine fits every stage of illness. A strengthening tonic harms a patient burning with high fever; ``Baptism and Penance are as purgative medicines, given to take away the fever of sin; whereas this sacrament is a medicine given to strengthen, and it ought not to be given except to them who are quit of sin.''

For a believer conscious of known grave sin, therefore, the Sacrament of Reconciliation is the ordinary remedy restoring the bond of charity before approaching Communion. Past venial sins do not bar reception---indeed, the Eucharist cleanses venial sins by arousing the fervor of charity---though clinging to actual venial sin at the moment of reception dulls the sacrament's actual spiritual refreshment (\emph{ST} III, q.\,79, a.\,8). And whenever sacramental reception is rightly deferred or physically impossible, the Church's distinction between \emph{sacramental eating} and \emph{spiritual eating} (\emph{ST} III, q.\,80, a.\,1, corpus and ad 3) ensures that the soul is not cut off from Christ: one may still unite oneself to him in faith, charity, and desire for the sacrament until worthy sacramental reception becomes possible, though actual reception produces the fruit of the sacrament more fully than desire alone.

\subsection*{Suffering Within Grace: Moral and Pastoral Clarity}

Christian teaching on suffering refuses two opposite falsehoods: treating affliction as a divine verdict proving secret guilt, and romanticizing pain as a good in itself to be courted, manufactured, or endured passively when relief is at hand. Five boundaries guard moral truth and pastoral sanity.

First, \textbf{suffering is an experienced evil, never good in itself and never presumptive punishment.} Bodily illness, bereavement, mental anguish, poverty, and injury are real privations of genuine created goods that God made good and wills creatures to flourish in. Neither Scripture nor Catholic doctrine permits looking at a suffering person---or at one's own affliction---and assuming that pain represents a direct tally of divine retribution.

Second, \textbf{removable suffering ought to be resisted and relieved.} Loving God never requires surrendering to preventable disease, ignoring competent medical or psychological care, tolerating injustice against oneself or the vulnerable, or seeking out bodily harm. St.\ Francis de Sales states the balanced Catholic rule in \emph{Introduction to a Devout Life}, Part III, Chapter III:

\begin{quote}
When any evil befalls you, apply the remedies that may be in your power, agreeably to the will of God; for to act otherwise would be to tempt divine Providence. Having done this, wait with resignation for the success it may please God to send; and, should the remedies overcome the evil, return him thanks with humility, but if, on the contrary, the evils overcome the remedies, bless him with patience.\endnote{Francis de Sales, \emph{Introduction to a Devout Life}, Part III, ch.\ III (Pustet undated, p.\ 129; confirmed against the French text of the 1832 P\'erisse Fr\`eres edition: ``Quand il vous arrivera du mal, apportez-y les rem\`edes qui vous seront possibles, et selon Dieu; car de faire autrement, ce seroit tenter la divine Providence'').}
\end{quote}

Refusing lawful remedies in the name of piety is not trust in Providence; Francis calls it \emph{tempting} Providence. Self-harm, quietism before remediable cruelty, and the imposition of generalized severe physical austerity as a universal badge of sanctity stand strictly outside authentic Christian charity.

Third, \textbf{when remedies are exhausted and unavoidable suffering remains, grace enables a free union with Christ's Passion---without adding anything to its redemptive worth.} Christ's sacrifice on Calvary is once-for-all, infinite, and completely sufficient for the redemption of the world. Nothing a creature suffers repairs a shortage in Christ's blood. Yet St.\ Paul writes from prison to the Colossians that he rejoices in his sufferings for their sake, to ``fill up those things that are wanting of the sufferings of Christ, in my flesh, for his body, which is the church'' (Col 1:24; cf.\ Rom 8:17--18; 2 Cor 4:17; 12:9). What remains to be completed is not the objective value of Christ's redemption, which is infinite, but its lived application in the members of his Body across history. Because the Head and the members form one Mystical Body, an affliction that cannot be removed need not remain bare waste; offered freely to the Father in union with Christ's offering, it becomes an act of intercession, solidarity, and love.\endnote{Pope Benedict XVI's \emph{Spe salvi} (2007), nos.\ 35--38, develops this same teaching: while every effort must be made to reduce suffering in the world, no technical or political effort can abolish all suffering; healing arrives through learning to suffer \emph{with} others and \emph{with} Christ, who entered human suffering from within.}

Fourth, \textbf{pain itself has zero merit.} Merit belongs strictly to a free human act proceeding from sanctifying grace and animated by charity. Aquinas states the exact boundaries of Christian merit in \emph{ST} I-II, q.\,114, aa.\,1, 4, and 5. Between an infinite Creator and a finite creature who owes everything to him, there can be no juridical justice of equal exchange as between commercial partners. Therefore, ``man's merit with God only exists on the presupposition of the Divine ordination'' (\emph{ST} I-II, q.\,114, a.\,1, corpus)---God freely ordains that free acts performed under grace attain the eternal reward he promises, making God debtor not to our leverage, but to his own faithful promise (ad 3).

Moreover, Aquinas demonstrates in \emph{ST} I-II, q.\,114, a.\,4 that merit rests chiefly with \emph{charity}, not with the intensity of toil or the sharpness of agony:

\begin{quote}
The act of faith is not meritorious unless ``faith ... worketh by charity'' (Gal.\ 5:6). So, too, the acts of patience and fortitude are not meritorious unless a man does them out of charity, according to 1 Cor.\ 13:3: ``If I should deliver my body to be burned, and have not charity, it profiteth me nothing.''\endnote{Aquinas, \emph{ST} I-II, q.\,114, a.\,4, ad 3; see also corpus and ad 2, clarifying that difficulty arising from reluctance actually \emph{lessens} merit, whereas charity heals reluctance and makes the will ready.}
\end{quote}

Fifth, \textbf{no human being can ever merit the first grace} (\emph{ST} I-II, q.\,114, a.\,5). Every meritorious deed presupposes grace as its root; before grace, human nature neither equals the supernatural order nor overcomes the barrier of sin. As Aquinas puts it (ad 2), ``God gives grace to none but to the worthy, not that they were previously worthy, but that by His grace He makes them worthy,'' adding from Job 14:4 that God alone makes the unclean clean.\endnote{Aquinas, \emph{ST} I-II, q.\,114, a.\,5, corpus and ad 2. The Council of Trent, Session VI, Decree on Justification, ch.\ 16, and the \emph{Catechism of the Catholic Church}, nos.\ 2007--2011, teach this exact framework: no one can merit the initial grace of conversion and justification; all subsequent Christian merit belongs primarily to God's grace and secondarily to the believer's free cooperation moved by the Holy Spirit in charity.}

\subsection*{Love of Neighbor: The Inseparable Test and Fruit}

If someone claims deep friendship with God, fervent Eucharistic devotion, and patient endurance, yet harbors contempt, indifference, or malice toward the person standing next to them, Scripture pronounces an uncompromising verdict:

\begin{quote}
If any man say: I love God, and hateth his brother; he is a liar. For he that loveth not his brother whom he seeth, how can he love God whom he seeth not? And this commandment we have from God, that he who loveth God love also his brother (1 John 4:20--21).
\end{quote}

Our Lord had already bound the two precepts into one indivisible frame: after naming love of God with whole heart, soul, and mind as the first commandment, he added immediately, ``And the second is like to this: Thou shalt love thy neighbour as thyself. On these two commandments dependeth the whole law and the prophets'' (Matt 22:39--40), and warned that mercy shown or withheld from the least of his brethren is shown or withheld from himself (Matt 25:40).

In \emph{ST} II-II, q.\,25, a.\,1, Aquinas shows why love of God and love of neighbor are not two separate virtues competing for finite room in the heart, like two separate accounts drawn from one shrinking purse. Every virtue is specified by its formal object---the reason under which it loves what it loves. Just as the eye sees light and sees color \emph{under the aspect of light} by one and the same act of vision, so charity loves God and loves the neighbor by one specifically identical habit:

\begin{quote}
Now the aspect under which our neighbor is to be loved, is God, since what we ought to love in our neighbor is that he may be in God. Hence it is clear that it is specifically the same act whereby we love God, and whereby we love our neighbor. Consequently the habit of charity extends not only to the love of God, but also to the love of our neighbor.\endnote{Aquinas, \emph{ST} II-II, q.\,25, a.\,1, corpus. In \emph{ST} II-II, q.\,23, a.\,1, ad 2--3, Aquinas explains why this same friendship extends even to sinners and enemies: when we deeply love a friend, for that friend's sake we love everyone who belongs to him, even if they hurt or oppose us.}
\end{quote}

That unity protects both loves at once. Loving our neighbor \emph{in God} neither turns the neighbor into an idol (as though a finite person were our last end, \emph{ST} II-II, q.\,25, a.\,1, ad 3) nor treats the neighbor as a mere instrument for scoring spiritual points. Rather, it wills for every human being---friend, stranger, and enemy alike---the highest good that exists: real justice, mercy, and bodily care here, ordered toward eternal life with God. And because the neighbor is visible, concrete, and often inconvenient, neighbor-love stands as the plain test that prevents purported love of God from evaporating into self-flattering solitude.

St.\ Th\'er\`ese of Lisieux grasped this ecclesial unity in the climactic discovery of her \emph{Story of a Soul} (Chapter XI). Feeling pulled toward every mission in the Church---apostle, prophet, doctor, martyr---she turned to St.\ Paul's hymn of charity in 1 Corinthians 12--13 and realized that charity is the living heart pumping blood to every limb:

\begin{quote}
O Jesus, my Love, at last I have found my vocation. My vocation is love! Yes, I have found my place in the bosom of the Church, and this place, O my God, Thou hast Thyself given to me: in the heart of the Church, my Mother, I will be LOVE!\endnote{Th\'er\`ese of Lisieux, \emph{Autobiography}, ch.\ XI (Taylor ed., p.\ 182).}
\end{quote}

Yet human love rarely begins at such clarity. Most believers discover that their love for God starts tangled with fear, need, and self-interest. How does a creature's frail affection actually grow toward that wholehearted self-gift without losing its own created identity? For that ascent---and for the rest in which it culminates---the tradition turns to St.\ Bernard of Clairvaux and St.\ Augustine.
